\begin{table}[htbp]
\centering
\caption{Cross-Sectional Entry Thresholds: Population by Current Store Count, Border vs.\ Interior}
\label{tab:entry_threshold}
\begin{tabular}{lcccc}
\toprule
Current store count $N$ & \multicolumn{2}{c}{Median population} & \multicolumn{2}{c}{Counties ($n$)} \\
& Interior & Border & Interior & Border \\
\midrule
0 & 20{,}536 & 15{,}488 & 66 & 65 \\
1 & 44{,}602 & 35{,}436 & 31 & 38 \\
2 & 56{,}526 & 59{,}352 & 29 & 20 \\
3 & 61{,}213 & 71{,}811 & 22 & 16 \\
4+ & 170{,}441 & 280{,}922 & 133 & 59 \\
\midrule
\multicolumn{5}{l}{Threshold ratio $S_2/S_1$: interior $=1.27$, border $=1.67$} \\
\bottomrule
\end{tabular}
\begin{minipage}{0.95\textwidth}
\vspace{0.3em}
\footnotesize
\textit{Notes:} Bresnahan-Reiss-style entry-threshold comparison \citep{BresnahanReiss1990,BresnahanReiss1991},
using the current (2026-08-05 snapshot) dispensary count for the nine-state entry-density universe. Border
counties support their first store at a smaller local population than interior counties (35{,}436 vs.\
44{,}602), consistent with a cross-border demand supplement to the local market. Border counties do
\emph{not} support a second store at a smaller population, however -- the opposite of what a simple
demand-supplement story predicts -- and require a larger relative population increase to go from one store
to two ($S_2/S_1=1.67$ vs.\ $1.27$ in interior counties). Cell sizes are modest (20 border and 29 interior
counties at $N=2$), so this comparison should be read as suggestive rather than decisive, but it does not
corroborate a shift in the equilibrium number of firms a market supports, even though the hazard results in
Table~\ref{tab:entry_hazard_by_rank} show a real effect on the \emph{speed} of reaching a second entrant. We
read the two results together as evidence that entry speed and the equilibrium entry threshold are
empirically dissociated in this setting, discussed further in Section~\ref{sec:results}.
\end{minipage}
\end{table}
